\section{Buổi 1}
\textbf{Câu 1: Giải thích về cách mạng công nghiệp chuyển đổi số và cho ví dụ}\\
Đây là giai đoạn mà công nghệ số được tích hợp sâu vào mọi lĩnh vực kinh tế và xã hội, làm thay đổi cách con người sản xuất, quản lý và sinh hoạt. Chuyển đổi số không chỉ là đưa máy tính hay Internet vào sử dụng, mà là:

\begin{itemize}
	\item Chuyển dữ liệu và quy trình sang dạng số
	
	\item Kết nối con người, máy móc và hệ thống
	
	\item Ra quyết định dựa trên dữ liệu và tự động hóa
\end{itemize}

Thay vì ngày xưa phải đi đến trực tiếp cửa hàng thì mới mua được hàng, phải đến trực tiếp quán ăn mới ăn được thì bây giờ mọi người có thể đặt giao hàng trực tuyến qua các ứng dụng, sau đó sẽ có các nhân viên giao đến tận nơi giúp mình.

Lúc trước, muốn thực hiện một giao dịch ngân hàng hay những hoạt động về tài chính thì phải đến trực tiếp ngân hàng thì mới thực hiện được. Nhưng bây giờ ta có thể thực hiện hoàn toàn trực tuyến, một cách nhanh chóng và tiện lợi thông qua các ứng dụng banking của ngân hàng.

\textbf{Câu 2: Đối với xe công nghệ hiện nay thì cách hoạt động của chúng chiếm ưu thế gì hơn so với tập đoàn Mai Linh, Vinasun lúc xưa?}

Khách hàng dễ dàng kiểm soát và theo dõi lộ trình của mình thông qua ứng dụng trên điện thoại của mình, có thể dự đoán thời gian đến và chia sẻ được hành trình cho người khác.

Đặt xe trực tuyến dễ dàng nhanh chóng trên ứng dụng của hãng xe thay vì phải tốn công sức so với việc gọi điện đến tổng đài như trước.

Rõ ràng, minh bạch về giá cước: mọi người có sẽ biết trước số tiền phải trả là bao nhiêu ngay trước khi xác nhận đặt xe. Thay vì như trước đây, taxi Mai Linh và Vinasun sẽ tính tiền dựa trên số km nên khách hàng sẽ không thể biết trước chính xác số tiền phải trả, gây ra nhiều sự bất tiện.

\textbf{Câu 3: Cho ví dụ về chuyển đổi số trong giáo dục}

Sử dụng giáo trình điện tử để thay thế tài liệu giấy, sách giáo khoa. Với bước tiến này thì sẽ giúp giảm chi phí để mua tài liệu, sách giáo khoa và tránh mang quá nhiều sách vở của nhiều môn trong một ngày cho học sinh, sinh viên. Ngoài ra, chúng còn giúp giáo viên có thể cập nhật nội dung giảng dạy một cách dễ dàng hơn.

Thay vì phải giảng dạy hay học tập trực tiếp thì bây giờ giáo viên và học sinh, sinh viên cũng có thể học tập, giảng dạy từ xa thông qua các ứng dụng trực tuyến như Zoom, Google Meet,... Điều này giúp giáo viên và học sinh tương tác, học tập cùng nhau một cách linh hoạt và hiệu quả.

\textbf{Câu 4: Triết lý của xe tự hành là gì? Tại sao vẫn đầu tư vào xe tự hành?}

Mục tiêu đầu tư vào xe tự hành là để tạo ra một hệ thống giao thông an toàn, hiệu quả, thông minh và tối ưu.

Ví dụ: Một ngày đi làm, giam xe tại nhà 8 tiếng. Về đến nhà thì để xe qua đêm lại giam xe 8 tiếng. => Tận dụng thời gian nhàn rỗi của xe ta có thể cho xe đi làm việc khác. Một ví dụ khác: Trong đợt dịch Covid vừa rồi, đối với một vài nơi cần người vận chuyển vật phẩm vào các khu vực ổ dịch, nếu để con người trực tiếp điều khiển xe đi vào trung tâm dịch thì rất nguy hiểm => Sử dụng xe tự hành thay thế để đảm bảo an toàn cho con người.

Tuy nhiên, vẫn còn nhiều thách thức nếu Việt Nam đưa xe tự hành vào sử dụng. Ví dụ đối với hệ thống giao thông chằng chịt như Việt Nam thì cần phải sửa đổi trước, phải xây dựng hay phân bố riêng làn cho xe tự hành. Nhưng những lợi ích tiềm năng của xe tự hành là rất lớn nên vẫn rất đáng để đầu tư.

\textbf{Câu 5: Cho ví dụ dữ liệu thị giác nằm ở đâu trong quá trình sản xuất và quy trình quản lý xã hội}

Dữ liệu thị giác nằm trong mọi lĩnh vực:

\begin{enumerate}
	\item Trong quản lý giao thông: Theo dõi tốc độ, phát hiện ùn tắc, đếm số lượng phương tiện giao thông, theo dõi người tham gia phương tiện di chuyển.
	
	\item Trong sản xuất thì thường dùng để quản lý các kho hàng như đếm số lượng hàng hóa, hay xác định vị trí sản phẩm khi cần tìm kiếm.
\end{enumerate}

\textbf{Câu 6: Nêu điểm khác biệt của mắt người và thiết bị số khi thu thập hình ảnh}

\begin{table}[H]
	\centering
	\begin{tabular}{|p{7cm}|p{7cm}|}
		\hline
		Mắt người & Thiết bị số \\ \hline
		Thu nhận ánh sáng bằng võng mạc & Thu nhận ánh sáng bằng cảm biến CCD hoặc CMOS \\ \hline
		Xử lý hình ảnh kết hợp sinh lý và nhận thức của não & Xử lý theo thuật toán \\ \hline
		Độ phân giải thay đổi theo vùng nhìn & Độ phân giải cố định theo số lượng pixel \\ \hline
		Vùng trung tâm nhìn rõ hơn vùng ngoại vi & Phân bố đồng đều trên toàn ảnh \\ \hline
		Thích nghi tốt với thay đổi ánh sáng & Phụ thuộc cảm biến và thuật toán xử lý \\ \hline
		Ghi nhớ chọn lọc và mang tính nhận thức & Dễ sao chép và xử lý \\ \hline
	\end{tabular}
\end{table}

\textbf{Câu 7: Hiện nay, con người giải quyết thách thức về lỗ hổng ngữ nghĩa như thế nào?}

10 năm trước thì đây là bài toán gian nan. Nhưng hiện nay nhờ deep learning, máy có thể học tốt để hiểu và dễ dàng phát hiện, phân lớp, truy vấn đối tượng rất tốt.

Ví dụ: Xe tự hành gắn nhiều sensor => Đề phòng trường hợp nếu sensor này bị hư do các yếu tố khách quan như thời tiết,.. thì vẫn còn cái khác.

\textbf{Câu 8: Phân biệt Computer Graphic, Image \& Video Processing và Computer Vision}

\begin{table}
	\centering
	\begin{tabular}{|c|p{4cm}|p{4cm}|p{4cm}|}
		\hline
		& Computer Graphics & Image \& Video Processing & Computer Vision \\ \hline
		Mục tiêu & Mô phỏng thế giới thực và thế giới ảo & Tăng cường chất lượng ảnh/video & Tăng cường khả năng hiểu ảnh/video cho máy tính \\ \hline
		Input & Điểm 2D, 3D, các loại đường cong & Ảnh/video & Ảnh/video \\ \hline
		Output & Các thực thể hình học 2D, 3D & Ảnh/video tăng cường & Ngữ nghĩa của ảnh/video \\ \hline
	\end{tabular}
\end{table}

\textbf{Câu 9: Cho ví dụ về AR và VR?}

Ví dụ về AR: Ứng dụng AR trong việc hướng dẫn đường đi trong bệnh viện giúp người dùng hay bệnh nhân tìm ra nơi họ cần đến ngay cả khi không có nhân viên bệnh viện. Giúp giảm thiểu thời gian, an toàn và dễ dàng quản lý mọi người.

Ví dụ về VR: Mô phỏng không gian phẫu thuật ảo. Các bác sĩ có thể sử dụng VR để thực hiện các ca phẫu thuật phức tạp trên cơ thể người ở trong môi trường ảo, giúp làm quen với quy trình phẫu thuật và nâng cao kỹ năng trước khi thực nghiệm trên cơ thể người.

\section{Buổi 2}

\textbf{Câu 1: Ý nghĩa của hệ thống thị giác thông minh đóng góp quy trình giám sát, quản lí thông minh như thế nào?}

Hệ thống thị giác thông minh đóng vai trò quan trọng trong quy trình giám sát và quản lý thông minh:

\begin{itemize}
	\item Tự động hóa quy trình giám sát, hệ thống có thể hoạt động 24/7 mà không cần sự can thiệp trực tiếp từ con người, giúp giảm thiểu sai sót do yếu tố con người gây ra. Phát hiện sự cố nhanh chóng các sự cố bất thường như cháy nổ, xâm nhập trái phép, hay các hoạt động không hợp lệ một cách nhanh chóng.
	
	\item Phân tích dữ liệu thời gian thực: Hệ thống cung cấp thông tin tức thì để hỗ trợ ra quyết định, như điều phối nhân lực, điều chỉnh quy trình sản xuất, hoặc tối ưu hóa hoạt động vận hành.
	
	\item Giúp giảm chi phí nhân công và tài nguyên, các máy móc hệ thống có thể tự động hóa các nhiệm vụ như kiểm tra chất lượng sản phẩm, theo dõi lưu lượng giao thông, hay giám sát môi trường thay thế con người.
	
	\item Giảm thiểu sai sót tính toán: Việc sử dụng AI và machine learning để liên tục cải thiện độ chính xác của các dự đoán và phát hiện.
	
	\item Giảm rủi ro tai nạn, hạn giúp con người hạn chế tiếp xúc với môi trường độc hại như khai thác mỏ, sản xuất hóa chất, xử lý rác thải độc hại,...
\end{itemize}

\textbf{Câu 2: Thế nào là thông minh?}

Thông minh là khả năng thích ứng của sự thay đổi hay nói một cách khác là khả năng nắm bắt bản chất của sự vật hiện tượng một cách nhanh chóng.

\textbf{Câu 3: Giải thích câu nói “Intelligence is the ability to adapt to change." - Stephen Hawking}

Câu nói "Intelligence is the ability to adapt to change." của Stephen Hawking nhấn mạnh rằng trí thông minh không chỉ đơn thuần là kiến thức hay kỹ năng cố định, mà còn là khả năng linh hoạt, thích ứng với những thay đổi không ngừng của môi trường xung quanh.

\textbf{Câu 4: Vì sao phát triển mảng học sâu là có đóng góp đột phá trong thị giác máy tính?}

Học sâu có khả năng tự động trích xuất và học các đặc trưng từ dữ liệu hình ảnh mà không cần phải thiết kế thủ công như các phương pháp truyền thống.

Các mạng học sâu có khả năng xử lý dữ liệu lớn và đa dạng, từ hình ảnh 2D, video, đến dữ liệu 3D. Điều này rất quan trọng vì dữ liệu thị giác máy tính thường có kích thước lớn và phức tạp.

Cho được kết quả vượt xa so với các phương pháp trước đây trong nhiều bài toán thị giác máy tính, như: Nhận diện đối tượng, phân đoạn ảnh, phát hiện đối tượng.

\textbf{Câu 5: Thách thức về lỗ hổng ngữ nghĩa của hệ thống thị giác là gì?}

Thách thức về lỗ hổng ngữ nghĩa (semantic gaps) trong hệ thống thị giác máy tính là một trong những thách thức lớn nhất mà các hệ thống thị giác máy tính phải đối mặt. Chẳng hạn như:

\begin{itemize}
	\item Hệ thống xử lý hình ảnh ở mức pixel, các giá trị số hoặc đặc trưng kỹ thuật, mà bản thân chúng không mang ý nghĩa ngữ nghĩa rõ ràng.
	
	\item Hệ thống có thể không hiểu được rõ hết được ngữ nghĩa mà con người muốn, chẳng hạn như "đây là một con mèo đang ngủ trên ghế".
	
	\item Ngoài ra ý nghĩa của một đối tượng trong hình ảnh còn phụ thuộc vào ngữ cảnh xung quanh. Ví dụ: Một con dao có thể là dụng cụ nhà bếp hoặc một vật nguy hiểm, tùy thuộc vào môi trường (nhà bếp hoặc một cảnh hành động).
	
	\item Một đối tượng hoặc cảnh có thể mang nhiều ý nghĩa khác nhau trong các tình huống khác nhau. Ví dụ: Một biểu cảm khuôn mặt có thể biểu thị niềm vui hoặc sự mỉa mai, tùy thuộc vào bối cảnh.
\end{itemize}

Nhưng hiện nay những thách thức về lỗ hổng ngữ nghĩa của hệ thống thị giác đã không còn là quá đáng quan ngại nữa vì Trí tuệ nhân tạo đã có khả năng giải quyết các thách thức đó.

\textbf{Câu 6: Hệ thống thị giác gồm các tầng nào? Xử lý ảnh và video học kiến thức ở tầng nào?}

Hệ thống thị giác gồm 3 tầng:

\begin{itemize}
	\item \textbf{Tầng 1:} Xử lý dữ liệu cơ bản
	
	\item \textbf{Tầng 2:} Xử lý các tác vụ đơn
	
	\item \textbf{Tầng 3:} Xử lý các ứng dụng phức hợp
\end{itemize}

Xử lý ảnh và video số học các kiến thức chủ yếu ở tầng 1 và 2.

\textbf{Câu 7: Cho một ứng dụng hệ thống thị giác ở tầng 3}

\textbf{Tên ứng dụng:} Hệ thống quản lý giao thông thông minh.

\textbf{Input:} Các video thô, ảnh thô.

\textbf{Output:}

\begin{itemize}
	\item \textbf{Data:} Ảnh đã được tăng cường.
	
	\item \textbf{Information:} Các cảnh quay có vật thể.
	
	\item \textbf{Knowledge:} Thông tin các xe, số lượng phương tiện, các làn đường.
	
	\item \textbf{Intelligence:} Điều khiển đèn giao thông ở mỗi lane, phân luồng giao thông. Xác định phương tiện vi phạm để xử phạt.
	
	\item \textbf{Wisdom:} Cảnh báo sớm tai nạn có thể xảy ra. Dự đoán tình trạng giao thông.
\end{itemize}

\section{Buổi 3}

\textbf{Câu 1: Màu trên vật đâu ra?}

Màu sắc trên một vật xuất hiện nhờ vào cách mà vật thể tương tác với ánh sáng. Khi ánh sáng chiếu vào bề mặt của một vật thể, vật đó sẽ hấp thụ các bước sóng nhất định và phản xạ các bước sóng khác. Các bước sóng phát ra sau đó truyền tới mắt chúng ta, giúp chúng ta cảm nhận được màu sắc.

\textbf{Câu 2: Tìm hiểu cách chuyển đổi ảnh màu sang ảnh độ sáng và ngược lại}

Ảnh màu được biểu diễn trong không gian màu, còn ảnh độ sáng chỉ chứa thông tin về độ sáng mà không có thông tin màu sắc. Quá trình chuyển đổi thường dựa vào việc lấy trung bình hoặc tỉ lệ giữa các kênh màu (R,G,B) và giá trị này sẽ là giá trị độ sáng của pixel tương ứng trong ảnh grayscale.
\[I_{gray} = (R + G + B) / 3\]

Mắt người nhạy cảm với ánh sáng xanh lá nhất (kênh G), sau đó là đỏ (R), và ít nhạy cảm với xanh lam (B) nên công thức chuẩn dựa trên mắt người là: $I_{gray} = 0.2989R + 0.5870G + 0.1140B$.

Chuyển đổi từ ảnh grayscale trở lại ảnh màu là một vấn đề phức tạp hơn, bởi vì ảnh độ sáng không chứa thông tin về màu sắc gốc. Một cách đơn giản là tạo 3 kênh màu đỏ, xanh lá và xanh dương, và gán giá trị của kênh grayscale cho cả 3 kênh: $R = G = B = I_{gray}$. Kết quả là một ảnh màu đơn sắc (grayscale nhưng ở dạng RGB).

\textbf{Câu 3: Vì sao 1 màu sắc trong ảnh số và video số thường được cấu thành từ 3 thành phần R, G, B?}

Nguyên do chính từ cấu tạo sinh học của mắt người và các nguyên lý vật lý về ánh sáng.

\begin{enumerate}
	\item \textbf{Cấu tạo của mắt người (Thuyết Tam sắc): }
	Mắt người có các tế bào thụ cảm ánh sáng gọi là tế bào nón (cones). Có 3 loại tế bào nón chính, mỗi loại nhạy cảm với một dải bước sóng ánh sáng khác nhau:
	
	\begin{itemize}
		\item \textbf{L-cones:} Nhạy với bước sóng dài (màu Đỏ).
		
		\item \textbf{M-cones:} Nhạy với bước sóng trung bình (màu Xanh lá).
		
		\item \textbf{S-cones:} Nhạy với bước sóng ngắn (màu Xanh dương).
	\end{itemize}
	
	Vì não bộ chúng ta tạo ra cảm giác về màu sắc bằng cách kết hợp tín hiệu từ 3 loại tế bào này, nên các kỹ sư chỉ cần sử dụng 3 nguồn sáng tương ứng (Red, Green, Blue) để "đánh lừa" mắt rằng chúng ta đang thấy hàng triệu màu sắc khác nhau.
	
	\item \textbf{Nguyên lý cộng màu (Additive Color Mixing): }
	Ảnh số chủ yếu được hiển thị trên các thiết bị tự phát sáng như màn hình điện thoại, máy tính hay TV. Các thiết bị này sử dụng mô hình màu RGB:
	
	\begin{itemize}
		\item Mỗi điểm ảnh (pixel) trên màn hình thực chất gồm 3 điểm màu siêu nhỏ: Đỏ, Xanh lá và Xanh dương.
		
		\item Khi trộn các luồng ánh sáng này với cường độ khác nhau, chúng tạo ra các màu mới. Ví dụ: Đỏ + Xanh lá = Vàng; Đỏ + Xanh lá + Xanh dương = Trắng.
		
		\item Với 3 thành phần này, chúng ta có thể tạo ra một dải màu (gamut) đủ rộng để đáp ứng hầu hết nhu cầu quan sát của con người.
	\end{itemize}
	
	\item \textbf{Hiệu quả về lưu trữ và xử lý: }
	Việc sử dụng 3 thành phần là một sự cân bằng hoàn hảo giữa chất lượng hình ảnh và dung lượng dữ liệu:
	
	\begin{itemize}
		\item \textbf{Đủ dùng:} 3 kênh màu đã đủ để tái tạo hình ảnh trông rất tự nhiên.
		
		\item \textbf{Tiết kiệm:} Nếu dùng ít hơn (ví dụ 2 màu), chúng ta không thể tạo ra đủ các sắc độ. Nếu dùng nhiều hơn (ví dụ 4-5 màu), dung lượng ảnh sẽ tăng lên rất nhiều trong khi mắt người khó có thể nhận ra sự khác biệt rõ rệt.
	\end{itemize}
	
	Tuy nhiên, trong in ấn, người ta lại dùng 4 thành phần (CMYK: Cyan, Magenta, Yellow, Black). Điều này là do mực in hoạt động theo nguyên lý trừ màu (hấp thụ ánh sáng) thay vì phát sáng, nên cần thêm màu đen (K) để tạo độ tương phản sâu mà 3 màu kia khi trộn lại khó đạt được một cách hoàn hảo.
\end{enumerate}

\textbf{Câu 4: Giải thích ý nghĩa của (H, S, V)}

Không gian màu HSV (Hue, Saturation, Value) là một hệ màu được sử dụng để tạo ra tất cả các màu thông qua ba thông số chính: Hue (H) biểu thị sắc màu, Saturation (S) thể hiện độ bão hòa của màu sắc, và Value (V) chỉ mức độ sáng của màu. Trong đó:

\begin{itemize}
	\item Hue xác định bản chất của màu sắc mà không phụ thuộc vào độ sáng hoặc độ bão hòa.
	
	\item Saturation càng cao, màu sắc càng rực rỡ và mạnh mẽ. Saturation thấp làm màu sắc trở nên nhợt nhạt, gần như chỉ còn là tông xám.
	
	\item Value quyết định xem màu sắc đó trông sáng, tối, hoặc gần như đen.
\end{itemize}

\textbf{Câu 5: Với $H = 18, S = 10, V = 10$, thực hiện Quantization, có thể biểu diễn được bao nhiêu màu, so sánh với số lượng màu mà không gian RGB có thể biểu diễn như thế nào? Vì sao có thể chấp nhận sử dụng xấp xỉ này?}

Với$ H = 18, S = 10, V = 10$, có thể biểu diễn được $18 * 10 * 10 = 1800$ màu, khi so sánh với so lượng màu trong không gian RGB (~16.8 triệu màu), xấp xỉ này vẫn có thể chấp nhận được. Do màu sắc biến thiên liên tục và không có sự chuyển tiếp rõ ràng. Trong khi đó mắt người cực kỳ nhạy với cường độ sáng (Value) nhưng không phân biệt được màu sắc (Hue), các kênh thần kinh dẫn truyền xử lý màu sắc và độ sáng tách biệt với nhau, phù hợp trực tiếp với ý tưởng hình thành không gian HSV.

\textbf{Câu 6: Chứng minh tại sao $||x||_1 = 1$ có quỹ tích là hình thoi?}
\begin{gather*}
	||x||_1 = 1 \Leftrightarrow |x_1| + |x_2| = 1 \\
	\Rightarrow (x_1 + x_2 = 1) \quad or \quad (x_1 - x_2 = 1) \quad or \quad (-x_1 + x_2 = 1) \quad or \quad (-x_1 - x_2 = 1)
\end{gather*}
Mặt khác, $|x_1|, |x_2 \le 1$ nên ta có quỹ tích là hình thoi.

\textbf{Câu 7: Chứng minh tại sao $||x||_{\infty} = 1$ nằm trên các cạnh hình vuông?}
\[
	||x||_{\infty} = 1 \Leftrightarrow max\{|x_1|, |x_2|\} = 1 \\
\]
Trường hợp 1:
\begin{gather*}
	|x_1| \le |x_2| \Rightarrow max\{|x_1|, |x_2|\} = |x_2| = 1 \\
	\Rightarrow x_2 = \pm 1 \quad and \quad |x_1| \le 1
\end{gather*}
Trường hợp 2:
\begin{gather*}
	|x_2| \le |x_1| \Rightarrow max\{|x_1|, |x_2|\} = |x_1| = 1 \\
	\Rightarrow x_1 = \pm 1 \quad and \quad |x_2| \le 1
\end{gather*}
Vậy ta có quỹ tích là hình vuông.

\textbf{Câu 8: Chứng minh tại sao $||x||_2 = 1$ nằm trên đường tròn?}
\begin{gather*}
	||x||_2 = 1 \Leftrightarrow \sqrt{(x_1^2 + x_2^2)} = 1 \\
	\Rightarrow x_1^2 + x_2^2 = 1
\end{gather*}
Vậy ta có quỹ tích là đường tròn.

\textbf{Câu 8: Khi nào ứng dụng các độ đo $||x||_1$, $||x||_2$, $||x||_{\infty}$ trong xử lý ảnh số và video số?}

\begin{itemize}
	\item \textbf{Độ đo Mahattan $||x||_1$}
	
	Khi dữ liệu nằm trên lưới, việc di chuyển bị giới hạn bởi các trục toạ độ (không được đi chéo).
	
	Ứng dụng: So sánh pixel trong kĩ thuật block matching, so sánh độ khác biệt giữ hai frame liên tiếp.
	
	\item \textbf{Độ đo Euclide $||x||_2$}
	
	Khi các chiều dữ liệu có tầm quan trọng tương đương nhau, khoảng cách cần tìm là khoảng cách ngắn nhất giữa hai điểm trong không gian thực mà không có vật cản.
	
	Ứng dụng: So sánh histogram màu, ứng dụng tính khoảng cách trong phân đoạn ảnh K-means, so sánh embedding trong nhận dạng khuôn mặt, ứng dụng trong lọc nhiễu Bilateral Filter, Non-local means.
	
	\item \textbf{Độ đo Chebyshev $||x||_{\infty}$}
	
	Khi các chiều dữ liệu có thể di chuyển độc lập, khoảng cách cần tìm chỉ phụ thuộc vào khoảng cách di chuyển xa nhất của các trục dữ liệu.
	
	Ứng dụng: Xác định vùng ảnh (bounding box, khoảng cách pixel trong lưới vuông), phát hiện biên cạnh, ứng dụng trong kĩ thuật dilation, erosion, phát hiện vùng chồng lấp (collision dectection), phát hiện cut scene.
\end{itemize}

\section{Buổi 4}

\textbf{Câu 1: Phân biệt nếu $H_t(i)$ chỉ lấy $N_t(i)$ thôi thì chúng khác gì với $\frac{N_t(i)}{N(t)}$}

Khi $H_t(i)$ chỉ lấy $N_t(i)$ thì kết quả sẽ phụ thuộc vào kích thước khung hình. Khung hình lớn hơn có thể dẫn đến giá trị $H_t(i)$ lớn hơn, ngay cả khi phân bố màu tương tự. Điều này dẫn đến nếu hai khung hình có kích thước khác nhau, giá trị $H_t(i)$ không phản ánh chính xác sự khác biệt về phân bố màu.

Ví dụ: Khung hình 1: Có 100 điểm ảnh, $N_t(i)=10$ (bin màu thứ $i$). Khung hình 2: Có 200 điểm ảnh, $N_t(i)=20$ (cùng bin màu).

\begin{itemize}
	\item Nếu dùng $H_t(i) = N_t(i)$:
	
	\begin{itemize}
		\item Khung hình 1: $H_t(i)=10$
		
		\item Khung hình 2: $H_t(i)=20$
	\end{itemize}
		
	Trông có vẻ khác nhau, nhưng thực tế phân phối màu là như nhau (đều chiếm 10\%).
		
	\item Nếu dùng $H_t(i) = \frac{N_t(i)}{N(t)}$:
	
	\begin{itemize}
		\item Khung hình 1: $H_t(i) = 10/100 = 0.1$
		
		\item Khung hình 2: $H_t(i) = 20/200 = 0.1$
	\end{itemize}
	
	Phân phối màu được chuẩn hóa, cho thấy cả hai khung hình có cùng tỉ lệ màu.
\end{itemize}

\textbf{Câu 2: Nếu ta không định lượng màu thì gây ra khó khăn gì?}

Khi không định lượng màu có thể dẫn đến một vài khó khăn như sau:

\begin{itemize}
	\item Không rõ ràng vì nếu không có cách định lượng, việc biểu diễn màu sắc chỉ dựa vào mô tả cảm quan (ví dụ: "đỏ nhạt", "xanh lá cây đậm") sẽ rất chủ quan và không đồng nhất. Và không thể so sánh chính xác được 2 màu có giống hoặc khác nhau không nếu nó khá tương đồng về màu sắc.
	
	\item Không thể tái tạo chính xác màu sắc đó nếu không có thông số định lượng vì các thiết bị như máy tính, máy ảnh, và màn hình yêu cầu giá trị định lượng cụ thể để hiển thị màu sắc. Dẫn đến việc tái tạo màu sắc chính xác là rất khó.
	
	\item Ảnh hưởng đến các lĩnh vực, ứng dụng thực tế khác như là công nghiệp(in, ấn, sơn nhà..), y tế,... Ví dụ như trong y tế thì phân tích ảnh X-quang hoặc MRI, phụ thuộc vào thông tin màu để phát hiện các bất thường. Nếu không có định lượng màu thì sẽ rất khó để chẩn đoán chính xác được.
\end{itemize}

\textbf{Câu 3: Màu là gì?}

Màu là một khái niệm trực quan và vật lý, được định nghĩa dựa trên sự tương tác giữa ánh sáng, vật thể, và hệ thống cảm nhận (như mắt người hoặc cảm biến).

Trong ngành thị giác máy tính, màu được định nghĩa là một đặc trưng trực quan được biểu diễn dưới dạng thông tin số, mô tả cách ánh sáng ở các bước sóng khác nhau được cảm nhận hoặc ghi nhận bởi hệ thống cảm biến (như camera) và được xử lý để phân tích hoặc tái tạo.

\textbf{Câu 4: Vân là gì?}

Vân là đặc trưng mô tả sự phân bố không gian của mức xám hoặc màu sắc, thể hiện tính lặp lại, độ thô, độ mịn và cấu trúc bề mặt của vùng ảnh. Vân giúp nhận dạng bản chất vật liệu, phân đoạn các vùng ảnh có cấu trúc giống nhau.

Trong thị giác máy tính, vân không phụ thuộc vào một pixel đơn lẻ mà phụ thuộc vào mối quan hệ giữa các pixel lân cận.

\textbf{Câu 5: Dáng là gì?}

Dáng là hình dạng hoặc cấu trúc tổng thể bên ngoài của một vật thể, giúp ta nhận diện và phân biệt nó với các vật thể khác. Dáng thường phản ánh các đặc điểm như đường nét, kích thước, và tỉ lệ, và nó có thể thay đổi tùy vào góc nhìn hoặc cách vật thể được sắp xếp.

Trong thị giác máy tính, dáng (hay hình dạng) được định nghĩa là tập hợp các điểm mà tại đó có biến thiên lớn về màu sắc.

\textbf{Câu 6: Nếu lược đồ bỏ đi các điểm ảnh không nằm trên biên cạnh thì có được hay không?}

Trong trường hợp xác định hình dáng của vật thể, việc loại bỏ các điểm ảnh không nằm trên biên cạnh có thể dẫn đến sai sót. Biên cạnh của vật thể trong ảnh là một yếu tố quan trọng để nhận diện và phân tách hình dáng, nhưng nếu loại bỏ các điểm ảnh không thuộc biên cạnh, ta có thể mất đi các thông tin quan trọng về cấu trúc tổng thể của vật thể, đặc biệt là khi đối mặt với các vật thể có biên cạnh không rõ ràng hoặc bị mờ. Hay nói một cách dễ hiểu là sẽ không thể hiện đầy đủ đặc trưng của ảnh.

\textbf{Câu 7: Khuyết điểm của lược đồ màu với lược đồ hệ số góc?}

\textbf{Color Histogram} có thể gặp khó khăn khi đối mặt với thay đổi ánh sáng, màu sắc giống nhau hoặc các điều kiện phức tạp về độ tương phản. Mặc dù hiệu quả trong các tác vụ nhận diện màu sắc, nhưng không phản ánh được cấu trúc và hình dáng của vật thể.

\textbf{Edge Direction Histogram} lại có ưu điểm lớn trong việc phân biệt hình dáng và cấu trúc biên của vật thể, đặc biệt trong các điều kiện thay đổi ánh sáng hoặc tương phản thấp. Tuy nhiên, EDH cũng có thể thiếu thông tin về màu sắc, điều này có thể làm giảm hiệu quả trong các tác vụ nhận dạng dựa trên màu sắc.

\textbf{Câu 8: Ứng dụng của lược đồ màu CH và lược đồ hệ số góc EDH}

\begin{enumerate}
	\item \textbf{Lược đồ màu:} Biểu diễn phân bố tần suất các mức màu trong ảnh, không phụ thuộc vị trí không gian của pixel.
	
	\begin{itemize}
		\item Truy vấn hình ảnh dựa theo nội dung (CBIR) bằng cách so sánh ảnh dựa trên phân bố màu.
		
		\item Phân loại ảnh theo chủ đề (biển, rừng, sa mạc, …).
		
		\item Phân đoạn ảnh dựa vào sự khác biệt màu sắc (K-means).
		
		\item Theo dõi đối tượng trong video bằng cách so sánh lược đồ màu của đối tượng giữa các frame.
		
		\item Phát hiện chuyển cảnh bằng cách so sánh lược đồ màu giữa các frame.
	\end{itemize}
	
	\item \textbf{Lược đồ hệ số góc:} Biểu diễn phân bố hướng của biên hoặc gradient trong ảnh.
	
	\begin{itemize}
		\item Nhận dạng đối tượng (người, xe, vật thể, …)
		
		\item Nhận dạng ký tự OCR
		
		\item Phân đoạn ảnh dựa trên cấu trúc
		
		\item Phân tích chuyển động trong video, phát hiện hướng chuyển động
	\end{itemize}
\end{enumerate}

\textbf{Câu 10: Vì sao với $g(x,y) = af(x,y)$ lại tăng độ tương phản màu sắc?}

Xét 2 điểm $(x, y)$ và $(x', y')$. Độ tương phản về giá trị độ xén giữa $(x, y)$ và $(x', y')$ trước và sau là:
\begin{gather*}
	Tp_t = |f(x, y) - f(x', y')| \\
	Tp_s = |g(x, y) - g(x', y')|
\end{gather*}
Xét mối quan hệ giữa $Tp_t$ và $Tp_s$:
\begin{align*}
	Tp_s &= |g(x, y) - g(x', y')| = |a.f(x, y) - a.f(x', y')| \\
		&= |a|.|f(x,y) - f(x',y')| = |a|.Tp_t
\end{align*}
Nếu $|a| \le 1$ thì $Tp_s \le Tp_t$, ngược lại thì $Tp_s > Tp_t$.

\textbf{Câu 11: Hàm định lượng $Q^{color}(I^{Color}(x, y))$ để tính mật độ các bin màu từ ảnh RGB}

Giả sử không gian HSV có \texttt{H = nH}, \texttt{S = nS}, \texttt{V = nV}, giá trị của H, S, V nằm trong đoạn $[0; 1]$. Mã giả:
\begin{lstlisting}
For each pixel in Image:
	Color_RGB = getColor(pixel)
	Color_HSV = RGB2HSV(Color_RGB)

	h = Color_HSV.h
	s = Color_HSV.s
	v = Color_HSV.v

	Index_h = (int)(h / (1/nH))
	Index_s = (int)(s / (1/nS))
	Index_v = (int)(v / (1/nV))
	Index = Index_h*nS*nV + Index_s*nV + Index_v

	H[Index] = H[Index] + 1
\end{lstlisting}

\section{Buổi 5}

\textbf{Câu 1: Làm sao chuyển toàn bộ về các ảnh có độ sáng biến thiên từ độ sáng biến thiên từ $[f_1, f_2]$ sang $[g_1,g_2]$?}

Gọi $f, g$ lần lượt là giá trị độ xám trước và sau khi thực hiện phép biến đổi tuyến tính $T(f) = af+b$. Ta có:
\begin{gather*}
	g_1 = a f_1 + b \\
	g_2 = a f_2 + b \\
	\Rightarrow g_2 - g_1 = a(f_2 - f_1) \\
	\Rightarrow a = \frac{g_2 - g_1}{f_2 - f_1} \\
	\Rightarrow b = g_1 - a f_1 = g_1 - \frac{g_2 - g_1}{f_2 - f_1}f_1 = \frac{g_1 f_2 - g_2 f_1}{f_2 - f_1}
\end{gather*}

Vậy phép biến đổi cần tìm là $g = T(f) = \frac{g_2 - g_1}{f_2 - f_1} f + \frac{g_1 f_2 - g_2 f_1}{f_2 - f_1}$.

\textbf{Câu 2: Chứng minh nếu $\Delta f_1 = \Delta f_2$ thì $\Delta g_1 > \Delta g_2$ cho hàm logarithm và $\Delta g_1 < \Delta g_2$ cho hàm mũ}

\textbf{Chứng minh $\Delta f_1 = \Delta f_2$ thì $\Delta g_1 > \Delta g_2$ cho hàm logarithm}

Đặt $\Delta f_1 = f_1’ - f_1, \Delta f_2 = f_2’ - f_2$

$\Rightarrow f_1’ = f_1 + \Delta f_1, f_2’ = f_2 + \Delta f_2$

Ta có:
\begin{gather*}
	\Delta g_1 = g_1' - g_1 = c*log(f_1') - c*log(f_1) = c*log(\frac{f_1'}{f_1}) = c * log(1 + \frac{\Delta f}{f_1})  \\
	\Delta g_2 = g_2' - g_2 = c*log(f_2') - c*log(f_2) = c*log(\frac{f_2'}{f_2}) = c * log(1 + \frac{\Delta f}{f_2})
\end{gather*}
Vì $logarithm$ là hàm đồng biến, $f_1 < f_2$

$\Rightarrow \Delta g_1 > \Delta g_2$

\textbf{Chứng minh $\Delta f_1 = \Delta f_2$ thì $\Delta g_1 < \Delta g_2$ cho hàm mũ}

Đặt $\Delta f_1 = f_1’ - f_1, \Delta f_2 = f_2’ - f_2$

$\Rightarrow f_1’ = f_1 + \Delta f_1, f_2’ = f_2 + \Delta f_2$

Ta có:
\begin{gather*}
	\Delta g_1 = g_1' - g_1 = e^{f_1'} - e^{f_1} = e^{f_1 + \Delta f_1} - e^{f_1} = e^{f_1}(e^{\Delta f} - 1) \\
	\Delta g_2 = g_2' - g_2 = e^{f_2'} - e^{f_2} = e^{f_2 + \Delta f_2} - e^{f_2} = e^{f_2}(e^{\Delta f} - 1)
\end{gather*}

Vì hàm mũ là hàm đồng biến, $f_1 < f_2$

$\Rightarrow \Delta g_1 < \Delta g_2$

\textbf{Câu 3: Lập bảng so sánh các phương pháp biến đổi màu dựa vào tuyến tính, phi tuyến, phân bố xác suất}

\begin{table}[H]
	\centering
	\begin{tabular}{|l|p{3cm}|p{4.5cm}|p{4.5cm}|}
		\hline
		\textbf{Tiêu chí} &
		\textbf{Tuyến tính} &
		\textbf{Phi tuyến} &
		\textbf{Phân bố xác suất} \\ \hline
		\textbf{Nguyên lý} &
		Dùng phép toán cộng/trừ hoặc ma trận để biến đổi pixel một cách tuyến tính. &
		Dùng các hàm phi tuyến (gamma, log, exp) để tăng cường độ sáng/tối cục bộ. &
		Tái phân phối giá trị pixel dựa trên histogram để cân bằng độ tương phản toàn ảnh. \\ \hline
		\textbf{Phương pháp} &
		Duyệt qua từng pixel rồi dùng phép biến đổi tuyến tính để thay đổi độ sáng hoặc độ tương phản. &
		Duyệt qua từng pixel rồi dùng phép biến đổi phi tuyến để thay đổi độ sáng hoặc độ tương phản. Tăng chi tiết vùng tối: $g(x,y) = c.log(f(x,y))$. Làm sáng mạnh vùng tối: $g(x,y) = e^{f(x,y)}$ &
		Tính histogram của ảnh. Tính CDF(I): $CDF(I) = \Sigma_{j=0}^I histogram(j)$. Tính giá trị pixel mới: $H'(I) = round(((CDF(I) - CDFmin) / (1 - CDFmin)) * (L-1))$ \\ \hline
		\textbf{Độ phức tạp} &
		Thấp: Chủ yếu thực hiện phép toán nhân, cộng. &
		Trung bình đến cao: Cần tính toán các hàm phi tuyến trên từng pixel. &
		Cao: Yêu cầu phân tích toàn bộ hình ảnh để tính toán phân bố xác suất và thực hiện các thay đổi. \\ \hline
		\textbf{Ưu điểm} &
		Nhanh, dễ triển khai, xử lý tốt trong các biến đổi đơn giản. &
		Tăng cường độ tương phản, điều chỉnh phù hợp với đặc tính mắt người, cải thiện chi tiết trong vùng tối/sáng. &
		Cân bằng cường độ ánh sáng toàn ảnh, tăng chi tiết ở cả vùng sáng và tối. \\ \hline
		\textbf{Nhược điểm} &
		Không hiệu quả trong việc tăng cường độ tương phản, có thể gây mất cân bằng màu. &
		Có thể gây méo màu hoặc mất chi tiết trong một số trường hợp, yêu cầu điều chỉnh tham số. &
		Phức tạp hơn, đòi hỏi tài nguyên tính toán lớn, có thể làm giảm tính tự nhiên của hình ảnh. \\ \hline
	\end{tabular}
\end{table}

\textbf{Câu 4: Phân biệt phương pháp scale ảnh và phương pháp super image resolution}

\begin{table}[H]
	\centering
	\begin{tabular}{|l|l|l|}
		\hline
		\textbf{Tiêu chí}   & \textbf{Scale ảnh}  & \textbf{Super image resolution}   \\ \hline
		Mục tiêu            & Thay đổi kích thước & Tăng độ phân giải (số pixel) thực \\ \hline
		Thông tin mới       & Không               & Có                                \\ \hline
		Nguyên lý           & Nội suy             & Học và tái tạo                    \\ \hline
		Chất lượng phóng to & Bị mờ đi            & Rõ nét                            \\ \hline
		Độ phức tạp         & Thấp                & Cao                               \\ \hline
		Phụ thuộc dữ liệu   & Không               & Có                                \\ \hline
	\end{tabular}
\end{table}

\section{Buổi 6}

\textbf{Câu 1: Trong toán tử Gaussian, tại sao những điểm gần tâm lại có trọng số lớn?}

Giá trị màu mới tại $(x,y)$ dùng toán tử Gaussian:
\begin{gather*}
	g(x,y) = \Sigma_i \Sigma_j f(x - i, y - j).h(i,j), \quad (i,j) \in O \\
	h(i, j) = \frac{1}{\srqt{2 \pi} \sigma} e^{-\frac{i^2 + j^2}{2 \sigma^2}}
\end{gather*}
Nên trọng số càng lớn thì mức độ ảnh hưởng (về giá trị màu) của điểm đó lên tâm càng lớn.

\textbf{Câu 2: Chứng minh công thức toán tử trung vị} 
\[g(x, y) = med\{f(x + i, y + j) | (i, j) \in O\}\]
thì
\[\Sigma |f(x + i, y + j) - med| = min \Sigma |f(x + i, y + j) - f(x + i', y + j')|\]

Giả sử ${f(x+i, y+j), \quad (i,j) \in O}$ được sắp thứ tự tăng dần và ký hiệu lại:
\[I_1 < I_2 < ... < I_n, \quad med = I_{n + 1}, \quad n=2v+1\]

Ta có:
\begin{align*}
	A &= \Sigma |f(x+i,y+j) - med| = \Sigma |f(x+i,y+j) - I_{v+1}| \\
	&= (I_{v+1} - I_1) + (I_{v+1} - I_2) + ... + (I_{v+1} - I_v) \\
	& + (I_{v+2} - I_{v+1}) + (I_{v+3} - I_{v+1}) + ... (I_{2v+1} - I_{v+1})\\
	&= (I_{v+2} + I_{v+3} + ... I_{2v+1}) - (I_1 + I_2 + ... I_v)
\end{align*}

Với $I_k$ bất kì thuộc $\{I_1, I_2, ..., I_{2v+1}\}$ xét:
\begin{align*}
	B &= \Sigma |f(x+i,y+j) - I_k| \\
	&= (I_k - I_1) + (I_k - I_2) + ... + (I_k - I_{k-1}) \\ 
	&+ (I_{k+1} - I_k) + (I_{k+2} - I_k) + ... (I_{2v+1} - I_k) \\
	&= (k - 1)I_k - (I_1 + I_2 + ... + I_{k-1}) - (2v - k + 1)I_k + (I_{k+1} + I_{k+2} + ... + I_{2v+1}) \\
	&= 2(k - 1 - v)I_k + (I_{k+1} + I_{k+2} + ... + I_{2v+1}) - (I_1 + I_2 + ... + I_{k-1})
\end{align*}

Cần chứng minh $A < B$ hay $A - B < 0$. Đặt $k = v + 1 + \alpha$, khi đó:
\begin{align*}
	B &= 2 \alpha I_k + (I_{k+1} + I_{k+2} + ... + I_{2v+1}) - (I_1 + I_2 + ... + I_{k-1}) \\
	A - B &= (I_{k+2} + ... + I_{v+2+\alpha - 1}) + (I_{v+1} + ... + I_{v+ \alpha - 1}) - 2 \alpha * I_{v + 1 + \alpha} < 0
\end{align*}
Do đó ta có điều phải chứng minh.

\textbf{Câu 3: Lập bảng thống kê 3 phương pháp làm trơn ảnh dựa trên miền không gian}

\begin{table}[H]
	\centering
	\begin{tabular}{|p{2cm}|p{4cm}|p{4cm}|p{4cm}|}
		\hline
		\multicolumn{1}{|c|}{\textbf{Tiêu chí}} &
		\multicolumn{1}{c|}{\textbf{Toán tử trung bình}} &
		\multicolumn{1}{c|}{\textbf{Toán tử Gaussian}} &
		\multicolumn{1}{c|}{\textbf{Toán tử trung vị}} \\ \hline
		\textbf{Nguyên lý} &
		Thay giá trị màu hiện thời bằng giá trị màu trung bình, giúp tránh xa giá trị min và max &
		Thay giá trị màu hiện thời phụ thuộc các điểm lân cận, trong đó các điểm gần tâm có sức ảnh hương lớn hơn, giúp tránh xa giá trị min và max. &
		Thay giá trị màu hiện thời bằng giá trị trung vị vì tổng độ dị biệt giữa trung vị và các giá trị khác trong lân cận là bé nhất. \\ \hline
		\textbf{Phương pháp} &
		1. Xác định kích thước cửa sổ k x k. 2. Tính trung bình các giá trị pixel trong cửa sổ. 3. Thay thế pixel trung tâm bằng giá trị trung bình. &
		1. Xác định kích thước cửa sổ k x k và độ lệch chuẩn σ. 2. Tạo mặt nạ (kernel) Gaussian. 3. Tích chập kernel với ảnh gốc. &
		1. Xác định kích thước cửa sổ k x k. 2. Lấy giá trị các pixel trong cửa sổ và sắp xếp theo thứ tự tăng dần. 3. Chọn giá trị trung vị (median) trong tập dữ liệu đó. \\ \hline
		\textbf{Khả năng làm mờ} &
		Làm mờ tổng quát, làm giảm nhiễu đều trên toàn ảnh nhưng không bảo toàn tốt chi tiết và cạnh. &
		Làm mờ tốt hơn toán tử trung bình, bảo toàn chi tiết hơn nhờ giảm trọng số ở vùng xa trung tâm. &
		Hiệu quả trong việc loại bỏ nhiễu muối tiêu mà không làm mờ quá mức chi tiết hoặc cạnh. \\ \hline
		\textbf{Ưu điểm} &
		Đơn giản, dễ triển khai. Hiệu quả với nhiễu Gaussian nhẹ. &
		Bảo toàn cạnh tốt hơn so với toán tử trung bình. Hiệu quả với nhiễu Gaussian. &
		Loại bỏ hiệu quả nhiễu muối tiêu. Giữ nguyên chi tiết và cạnh tốt. \\ \hline
		\textbf{Nhược điểm} &
		Làm mờ cạnh và mất chi tiết. Không hiệu quả với nhiễu dạng muối tiêu. &
		Phức tạp hơn toán tử trung bình. Không hiệu quả với nhiễu muối tiêu. &
		Ít hiệu quả với nhiễu Gaussian. Phức tạp hơn toán tử trung bình khi triển khai. \\ \hline
	\end{tabular}
\end{table}

\textbf{Câu 4: Vì sao toán tử trung bình có trọng số gradient lại khắc phục được làm mờ biên cạnh}

Ta cần xét giá trị màu mới tại biên cạnh. Xét $(m,n)$ ở biên:
\begin{itemize}
	\item Nếu $(i,j)$ không gần biên: $|f(m,n) - f(i,j)|$ lớn $\Rightarrow \delta = \frac{1}{|f(m,n) - f(i,j)|}$ nhỏ $\Rightarrow h(i,j)$ nhỏ (trọng số nhỏ)
	
	\item Nếu $(i,j)$ ở gần biên hoặc nằm trên biên: $|f(m,n) - f(i,j)|$ nhỏ $\Rightarrow \delta = \frac{1}{|f(m,n) - f(i,j)|}$ lớn $\Rightarrow h(i,j)$ lớn (trọng số lớn)
\end{itemize}
Giá trị màu mới của biên bị ảnh hưởng nhiều bởi các điểm xung quanh, ít bị ảnh hưởng bởi các điểm xa nó. Do đó biên không bị mờ đáng kể.

\textbf{Câu 5: Toán tử trung bình dựa trên mặt nạ quay khác gì với toán tự trung vị?}

\begin{itemize}
	\item \textbf{Toán tử trung bình:} Chỉ lấy 1 mask
	
	\item \textbf{Toán tử trung bình dựa trên mặt nạ quay:} lấy nhiều mask $\Rightarrow$ chắc chắn tốt hơn vì sử dụng nhiều dữ liệu hơn
\end{itemize}

\textbf{Câu 6: Giải thích tại sao trong công thức độ lệch chuẩn cần bình phương?}

\begin{itemize}
	\item Với $|x| > 1: x^2 > x ⇒$ Quan tâm những giá trị có ảnh hưởng lớn (trọng số cao)

	\item Với $|x| < 1: x^2 < x ⇒$ Ít quan tâm những điểm có ảnh hưởng nhỏ (trọng số nhỏ)
\end{itemize}

\section{Buổi 7}

\textbf{Câu 1: Tìm biên cạnh theo đạo hàm thì có khuyết điểm gì?}

Thực tế hay gặp trường hợp biên độ nhỏ (f' rất bé), gradient do đó cũng rất bé, ta cần phải lựa chọn một ngưỡng phù hợp để có thể phân biệt biên và điểm lân cận.

Giải pháp: Lấy f'' rồi tìm khoảng f'' đổi dấu (Zero crossing).

\textbf{Câu 2: Khuyết điểm của các phương pháp ở trong toán tử Gradient là gì?}

Phát hiện biên cạnh nhưng đồng thời các điểm nhiễu cũng được cho vào.

\textbf{Câu 3: Vì sao Prewitt thì hy vọng chống chịu nhiễu tốt hơn?}

Vì lấy nhiều cặp điểm (cụ thể là 3 cặp) và nhờ vào trung bình cộng để kéo giá trị xuống.

\textbf{Câu 4: Tại sao phải dùng phương pháp tích chập?}

Vì viết dạng khai triển tuy rõ ràng nhưng dài. Viết tích chập chỉ cần viết ma trận ngắn gọn mà vẫn đầy đủ.

\textbf{Câu 5: Chứng mình toán tử Laplace bất biến với phép quay}

Xét phép quay hệ tọa độ với góc quay $\theta$, ta có biến đổi tọa độ:
\begin{align*}
	x' = x.cos(\theta) + y.sin(\theta) \\
	y' = -x.sin(\theta) + y.cos(\theta)
\end{align*}

Ta cần biểu diễn các đạo hàm theo biến $(x',y')$ để chứng minh rằng:
\[\nabla'^2f = \frac{\partial ^2f}{\partial x'^2} + \frac{\partial ^2f}{\partial y'^2} = \nabla^2f\]
Thật vậy, theo quy tắc đạo hàm chuỗi, ta có:
\begin{align*}
	\frac{\partial f}{\partial x'} = \frac{\partial f}{\partial x} \frac{\partial x}{\partial x'} + \frac{\partial f}{\partial y} \frac{\partial y}{\partial x'} = cos(\theta) \frac{\partial f}{\partial x} + sin(\theta) \frac{\partial f}{\partial y}  \\
	\frac{\partial f}{\partial y'} = \frac{\partial f}{\partial x} \frac{\partial x}{\partial y'} + \frac{\partial f}{\partial y} \frac{\partial y}{\partial y'} = -sin(\theta)\frac{\partial f}{\partial x} + cos(\theta) \frac{\partial f}{\partial y} \\
\end{align*}
Đạo hàm bậc 2 ta có:
\begin{align*}
	\frac{\partial ^2 f}{\partial x'^2} = \left(cos(\theta) \frac{\partial f}{\partial x} + sin(\theta) \frac{\partial f}{\partial y}\right)^2 \\
	\frac{\partial ^2 f}{\partial y'^2} = \left(-sin(\theta)\frac{\partial f}{\partial x} + cos(\theta) \frac{\partial f}{\partial y}\right)^2
\end{align*}
Cộng vế theo vế, ta có:
\[\frac{\partial ^2 f}{\partial x'^2} + \frac{\partial ^2 f}{\partial y'^2} = (cos^2(\theta) + sin^2(\theta)) (\frac{\partial ^2 f}{\partial x^2} + \frac{\partial ^2 f}{\partial y^2}) = \frac{\partial ^2 f}{\partial x^2} + \frac{\partial ^2 f}{\partial y'^2}\] (đpcm)

\textbf{Câu 6: Chứng minh: $(\nabla ^2 f)(x, y) \approx f(x + 1, y) + f(x - 1, y) + f(x, y + 1) + f(x, y - 1) - 4f(x, y)$}

Ta có:
\[(\nabla ^ 2 f)(x, y) = \frac{\partial ^2 f}{\partial x^2} + \frac{\partial ^2 f}{\partial y^2}\]

Cần tính $\frac{\partial ^2 f}{\partial x^2}$ và $\frac{\partial ^2 f}{\partial y^2}$

Có: $f_x(x, y) \approx f(x + 1, y) - f(x, y)$

Áp dụng điều này:

\[\frac{\partial ^2 f}{\partial x^2} \approx f_x(x + 1, y) - f_x(x, y)\]

Lại có:
\begin{gather*}
	f_x(x + 1, y) \approx f(x+1, y) - f(x, y) \\
	f_x(x, y) \approx f(x, y) - f(x-1, y) \\
	\Rightarrow \frac{\partial ^2 f}{\partial x^2} \approx f(x+1, y) + f(x-1, y) - 2f(x, y)
\end{gather*}

Tương tự:
\[\frac{\partial ^2 f}{\partial y^2} \approx f(x, y+1) + f(x, y-1) - 2f(x, y)\]

Từ đó:
\[\frac{\partial ^2 f}{\partial x^2} + \frac{\partial ^2 f}{\partial y^2} \approx f(x + 1, y) + f(x - 1, y) + f(x, y + 1) + f(x, y - 1) - 4f(x, y)\]

\section{Buổi 8}

\textbf{Câu 1: Chứng minh: $f(x, y) = (1 - a)(1 - b).f(l, k) + a(1 - b).f(l + 1, k) + b(1 - a).f(l, k + 1) + ab.f(l + 1, k + 1)$}

Với $l = round(x), a = x - l, k = round(y), b = y - k$
Dựa trên phương trình đoạn thẳng đi qua 2 điểm $A(x_0, y_0)$ và $B(x_1, y_1)$:
\[y = (1 - t)y_0 + ty_1, \quad t \in [0; 1]\]

Ta có:

\[f(x, y) = b.f(x, k + 1) + (1 - b).f(x, k) \quad (*)\]

Lại có:
\begin{align*}
	f(x, k + 1) &= a.f(l + 1, k + 1) + (1 - a).f(l, k + 1) \\
	f(x, k) &= (1 - a).f(l, k) + a.f(l + 1, k)
\end{align*}

Thế vào (*):

$\Rightarrow f(x,y) = (1-a)(1-b).f(l,k) + a(1-b).f(l+1,k) + b(1-a).f(l,k+1) + ab.f(l+1,k+1)$

\textbf{Câu 2: Phân biệt phân đoạn ảnh và phân lớp ảnh?}

Phân lớp là quá trình gán một class ID cho toàn bộ ảnh (Single Class Classification) hoặc nhiều class ID cho toàn bộ ảnh (Multiclass Classification).

Phân đoạn là quá trình gán ID cho từng điểm ảnh trong một ảnh, tương đương bài toán Classification cho điểm ảnh thay vì toàn bộ ảnh.

\textbf{Câu 3: Phương pháp Elbow để xác định số cụm k trong K-means clustering}

Phương pháp Elbow dùng để chọn số cụm k tối ưu trong k-means bằng cách quan sát mức độ giảm của sai số trong cụm khi tăng k.

\begin{itemize}
	\item Khi k tăng, các điểm dữ liệu được chia nhỏ hơn
	
	\item Tổng sai số trong cụm sẽ giảm dần
	
	\item Tuy nhiên, sau một giá trị k nào đó, mức giảm trở nên không đáng kể
\end{itemize}}

Điểm mà tại đó tốc độ giảm thay đổi rõ rệt được gọi là điểm khuỷu tay (elbow).

Thường dùng WCSS (Within-Cluster Sum of Squares), còn gọi là SSE:
\[WCSS = \Sigma_{i=1}^k \Sigma_{x \in C_i} ||x - \mu_i||^2\]
Trong đó:

\begin{itemize}
	\item $C_i$ là cụm thứ $i$.
	
	\item $\mu_i$ là tâm cụm.
	
	\item $||x - \mu_i||^2$ là bình phương khoảng cách Euclide.
\end{itemize}

Thuật toán:

\begin{enumerate}
	\item Cho k chạy từ 1 đến Kmax
	
	\item Chạy thuật toán K-means với từng giá trị k
	
	\item Tính WCSS tương ứng cho mỗi k
	
	\item Vẽ đồ thị WCSS theo k để xác định elbow: Trước elbow, tăng k làm WCSS giảm mạnh. Sau elbow, tăng k giảm WCSS rất ít. Giá trị k tại elbow được chọn là k tối ưu.
\end{enumerate}

\textbf{Câu 4: Phương pháp Silhoulette để xác định số cụm k trong K-means clustering}

Phương pháp Silhouette dùng để đánh giá chất lượng phân cụm và xác định số cụm k tối ưu dựa trên mức độ:

\begin{itemize}
	\item Gắn kết trong cụm
	
	\item Tách biệt giữa các cụm
\end{itemize}

Mỗi điểm dữ liệu được gán một hệ số Silhouette nằm trong khoảng từ -1 đến 1.

Với một điểm dữ liệu $x$:

\begin{itemize}
	\item $a(x)$: khoảng cách trung bình từ $x$ đến các điểm khác trong cùng cụm.
	
	\item $b(x)$: khoảng cách trung bình nhỏ nhất từ $x$ đến các điểm trong cụm khác gần nhất
\end{itemize}

Hệ số Silhouette của điểm $x$ được tính như sau:
\[s(x) = \frac{b(x) - a(x)}{max\{a(x), b(x)\}}\]
\begin{itemize}
	\item $s(x) \approx 1$: điểm được phân cụm rất tốt.
	
	\item $s(x) \approx 0$: điểm nằm gần ranh giới giữa các cụm.
	
	\item $s(x) < 0$: điểm có khả năng bị gán sai cụm.
\end{itemize}

Giá trị Silhouette trung bình của toàn bộ dữ liệu dùng để đánh giá chất lượng phân cụm tổng thể.

Thuật toán:
\begin{enumerate}
	\item Cho k chạy từ 1 đến Kmax
	
	\item Chạy thuật toán K-means với từng giá trị k
	
	\item Tính hệ số Silhouette cho từng điểm
	
	\item Tính Silhouette trung bình của toàn bộ dữ liệu
	
	\item Chọn k sao cho Silhouette trung bình là lớn nhất
\end{enumerate}

\section{Buổi 9}

\textbf{Câu 1: Ứng dụng của 2 công thức sau là gì?}
\begin{gather*}
	\zeta \{(f * h)(x, y)\} = F(u, v) . H(u, v) \\
	\zeta \{f(x, y).h(x, y)\} = (F * H)(u, v)
\end{gather*}

Định lý tích chập cho phép thay thế phép tích chập phức tạp trong miền không gian bằng phép nhân đơn giản trong miền tần số, giúp việc tính toán đơn giản và hiệu quả hơn. Nhờ có FFT, việc lọc tín hiệu và ảnh trở nên hiệu quả hơn.

\textbf{Câu 2: Thông số $D_0$ ảnh hưởng như thế nào đối với bộ lọc tần số?}

Đối với bộ lọc thông thấp (Lowpass), thường dùng để khử nhiễu, khi giá trị $D_0$ tăng, chỉ những thành phần nhiễu có tần số cao hơn mới bị loại bỏ, do đó nhiễu phải đủ mạnh hoặc đủ cao tần thì bộ lọc mới có tác dụng rõ rệt.

Đối với bộ lọc thông cao (Highpass), thường dùng để làm nổi bật biên cạnh, khi giá trị $D_0$ tăng, chỉ các biên có thành phần tần số rất cao mới được giữ lại, còn các biên yếu hoặc thấp tần sẽ bị suy giảm.

$\Rightarrow$ Khi $D_0$ tăng, bộ lọc thông thấp trở nên ít khắt khe hơn trong việc loại bỏ nhiễu, còn bộ lọc thông cao trở nên khắt khe hơn trong việc giữ lại biên cạnh.
























